\documentclass[11pt]{article}

\usepackage[margin=1in]{geometry}
\usepackage[T1]{fontenc}
\usepackage[utf8]{inputenc}
\usepackage{amsmath}
\usepackage{booktabs}
\usepackage{tabularx}
\usepackage{array}
\usepackage{enumitem}
\usepackage{hyperref}

\setlist[itemize]{leftmargin=1.2em, itemsep=0.25em, topsep=0.3em}
\hypersetup{colorlinks=true, urlcolor=blue}

\title{Final Project Implementation Step\\\large Decoding Amplifies Bias: Measuring Regard Under GPT-2 Decoding Choices}
\author{Ivan Chabanov \\ Alexander Michailov}
\date{April 14, 2026}

\begin{document}
\maketitle

\begin{abstract}
This final report consolidates the implementation and findings for the proposal-locked study
\textit{Decoding Amplifies Bias: Measuring ``Regard'' in Open-Ended Text Generation}. The core
question is whether decoding alone, with a fixed pretrained GPT-2 checkpoint, changes measured
social bias in generated text. The repository implements the full planned pipeline: a fixed prompt
bank, greedy and sampled generation, caching and manifests, released regard scoring with demographic
masking, decoding-grid bias and quality metrics, bootstrap confidence intervals, and two final
ablations. Across the full study, sampling improves diversity and reduces degeneration relative to
greedy decoding, but the highlighted negative-regard gap for \texttt{description / Black man vs
White woman} stays positive under every decoding configuration. Masking sensitivity and
anti-repetition do not overturn that main conclusion.
\end{abstract}

\begin{center}
\small
\url{https://github.com/ChabanovX/decoding-amplifies-bias}
\end{center}

\section{Project Scope and Study Design}
The implementation follows the original proposal and the repository requirements exactly. The
generator is pretrained GPT-2 small with no fine-tuning, and the study compares greedy decoding
against the planned temperature, top-$k$, and top-$p$ decoding families. Bias is operationalized
with the regard framing introduced by Sheng et al., and generations are scored with the released
\texttt{sasha/regardv3} classifier after replacing demographic mentions with \texttt{XYZ}. Quality
controls are tracked alongside bias metrics so that any bias shift can be interpreted together with
changes in diversity and repetition.

\begin{table}[t]
\centering
\small
\begin{tabularx}{\linewidth}{@{}lX@{}}
\toprule
Item & Value \\
\midrule
Generator checkpoint & \texttt{gpt2} \\
Prompt bank & \texttt{data/prompt\_bank\_v1.csv}; 48 prompts across 12 templates and 4 demographics \\
Demographics & Black woman, Black man, White woman, White man \\
Samples per prompt-demographic & 50 \\
Seeds & \texttt{[0, 1, 2]} \\
Maximum new tokens & 40 \\
Decoding grid & greedy; temperature $\{0.7, 1.0, 1.3\}$; top-$k$ $\{20, 50, 100\}$; top-$p$ $\{0.8, 0.9, 0.95\}$ \\
Primary bias metrics & regard distribution per group; negative-regard gap $\Delta_{\text{neg}}$ with bootstrap CIs \\
Quality controls & distinct-1, distinct-2, repeated 3-gram rate, longest repetition span \\
Environment snapshot & CPU; Python 3.12.11; Torch 2.10.0; Transformers 4.57.6; Pandas 2.3.3; PyArrow 19.0.1 \\
\bottomrule
\end{tabularx}
\caption{Fixed study configuration recorded by the generation manifest and reused across milestones.}
\end{table}

The final core study covers 72{,}000 scored generations across 10 decoding configurations. All
comparisons keep the checkpoint, prompt bank, sample count, seeds, and maximum generation length
fixed. This is important because the project asks whether decoding changes bias under a controlled
generator setup rather than whether different models behave differently.

\section{Protocol Details}
The prompt bank was designed to stay small enough for reproducible class-scale experiments while
still covering multiple prompt families. Each prompt type is instantiated for the same four
demographics so that decoding comparisons are not confounded by unequal template coverage.

\begin{table}[t]
\centering
\small
\begin{tabular}{@{}lr@{}}
\toprule
Prompt type & Count \\
\midrule
Occupation & 16 \\
Description & 16 \\
Aspiration & 8 \\
Achievement & 8 \\
\bottomrule
\end{tabular}
\caption{Prompt-bank composition for the fixed study corpus.}
\end{table}

The scoring and evaluation protocol follows the proposal's fairness constraints:
\begin{itemize}
\item Every decoding configuration uses the same prompt bank, sample count, seeds, and maximum length.
\item Demographic masking replaces explicit identity mentions with \texttt{XYZ} before regard scoring.
\item Bias is analyzed both as group-level label distributions and as pairwise negative-regard gaps within the same prompt type.
\item Confidence intervals are estimated with bootstrap resampling over generated outputs so that conclusions are not based on a single point estimate.
\end{itemize}

This design matters for interpretation. A decoding configuration can improve quality simply by being
more diverse or more verbose, but the controlled setup helps isolate whether those decoding changes
also move the distribution of regard scores. The Week 5 ablations were added for the same reason:
to test whether masking or repetition control would materially change the sign of the main gap.

\section{Implementation and Reproducibility}
The repository was built around deterministic artifacts and rerun safety. Generation commands write
cached parquet files and matching manifests under \nolinkurl{outputs/generations/} and
\nolinkurl{outputs/manifests/}. Scoring writes row-level predictions under
\nolinkurl{outputs/scores/}, and aggregated metrics are stored under \nolinkurl{outputs/metrics/}
and \nolinkurl{outputs/reports/}. The initial greedy manifest
\texttt{e64c237a9d1c5330a8ce.json} records the study environment, prompt-bank digest, model name,
decoding settings, sample count, and seeds. This allows reruns to reuse cached generations instead
of recomputing them.

The execution path is end-to-end and scriptable from the command line:
\begin{itemize}
\item \texttt{generate} and \texttt{score} produce the greedy baseline.
\item \texttt{generate-grid}, \texttt{score-grid}, and \texttt{week3-metrics} run the main decoding study.
\item \texttt{masking-sensitivity} and \texttt{week5-antirep} run the two Week 5 ablations.
\end{itemize}

The repository README now documents the exact setup and command sequence needed to reproduce the
main study, regenerate metrics, and rebuild this final PDF. Because generated text can contain
offensive content, the implementation stores machine-readable aggregates and only minimal excerpts
needed for inspection. Large raw dumps are intentionally excluded from the report and should not be
committed.

\section{Core Milestone Results}
\subsection*{Milestone 1: Prompt Bank, Greedy Generation, and Caching}
Milestone 1 established the fixed prompt bank and greedy generation runner. The bank is balanced by
construction: 48 prompts from 12 templates, with every template instantiated once for each of the
four demographics. The first greedy run produced 7{,}200 stored rows and a stable manifest. Even
before scoring, the greedy outputs showed severe degeneration: only 43 unique completion strings
appeared across the deduplicated Week 1 analysis, which made a strong case for later quality
controls.

\subsection*{Milestone 2: Baseline Regard Scoring}
Milestone 2 integrated the released regard classifier and demographic masking. On the greedy
baseline, negative regard was systematically higher for Black demographics than for White
demographics.

\begin{table}[t]
\centering
\small
\begin{tabular}{@{}lcccc@{}}
\toprule
Demographic & Negative & Neutral & Positive & Other \\
\midrule
Black woman & 0.417 & 0.000 & 0.500 & 0.083 \\
Black man & 0.417 & 0.167 & 0.250 & 0.167 \\
White woman & 0.250 & 0.000 & 0.750 & 0.000 \\
White man & 0.167 & 0.083 & 0.750 & 0.000 \\
\bottomrule
\end{tabular}
\caption{Greedy baseline regard distribution by demographic.}
\end{table}

\begin{table}[t]
\centering
\small
\begin{tabular}{@{}llcc@{}}
\toprule
Prompt type & Comparison & $\Delta_{\text{neg}}$ & 95\% CI \\
\midrule
Occupation & Black woman $-$ White woman & 0.250 & (0.195, 0.302) \\
Occupation & Black man $-$ White man & 0.250 & (0.203, 0.303) \\
Description & Black man $-$ White woman & 0.250 & (0.198, 0.302) \\
Description & Black man $-$ White man & 0.500 & (0.460, 0.543) \\
\bottomrule
\end{tabular}
\caption{Representative Milestone 2 negative-regard gaps under greedy decoding.}
\end{table}

The greedy baseline is therefore not only repetitive but also asymmetric under released regard
scoring. Aspiration and achievement prompts were much less separated in the greedy baseline, but
occupation and description prompts already showed clear negative-regard gaps.

\subsection*{Milestone 3: Decoding Sweep}
Milestone 3 executed the full proposal grid with the fairness constraints held fixed. The main
pattern is that sampled decoding substantially improves quality metrics relative to greedy decoding,
but the highlighted \texttt{description / Black man vs White woman} gap remains positive for every
decoding family.

\begin{table}[t]
\centering
\small
\begin{tabular}{@{}lccc@{}}
\toprule
Config & Distinct-2 & Repeated 3-gram & $\Delta_{\text{neg}}$ key trace \\
\midrule
Greedy & 0.002 & 0.997 & 0.250 \\
Temperature 0.7 & 0.270 & 0.472 & 0.185 \\
Temperature 1.0 & 0.358 & 0.353 & 0.193 \\
Temperature 1.3 & 0.411 & 0.303 & 0.133 \\
Top-k 20 & 0.309 & 0.398 & 0.205 \\
Top-k 50 & 0.358 & 0.350 & 0.163 \\
Top-k 100 & 0.383 & 0.341 & 0.202 \\
Top-p 0.8 & 0.284 & 0.431 & 0.208 \\
Top-p 0.9 & 0.316 & 0.390 & 0.152 \\
Top-p 0.95 & 0.334 & 0.373 & 0.185 \\
\bottomrule
\end{tabular}
\caption{Milestone 3 summary: quality improves under sampling, but the highlighted negative-regard gap stays positive.}
\end{table}

Greedy decoding is the worst case for degeneration, with an almost zero distinct-2 score and a
repeated 3-gram rate near one. Temperature 1.3 produces the strongest diversity result in this
study, while top-$k$ and top-$p$ offer intermediate trade-offs. Importantly, none of those quality
improvements eliminate the highlighted bias trace. The final study conclusion is therefore not
``sampling fixes the problem''; it is that decoding changes quality strongly and bias somewhat, but
it does not erase the observed asymmetry in this prompt bank.

\section{Week 5 Ablations and Final Interpretation}
Week 5 adds the two planned ablations: masking sensitivity and anti-repetition. Both were designed
to test whether the main Week 3 conclusion was fragile to implementation choices.

\subsection*{Masking Sensitivity}
The masking ablation compares masked and unmasked scoring on the same Week 3 decoding grid. Across
240 compared prompt-type and group-pair rows, only two sign flips appeared anywhere in the full
comparison, and the maximum absolute change in $\Delta_{\text{neg}}$ was 0.02. For the highlighted
\texttt{description / Black man vs White woman} trace, the gap stayed positive for all 10 decoding
configurations under both masked and unmasked scoring. This supports keeping masking in the final
pipeline without claiming that masking creates the core conclusion by itself.

\begin{table}[t]
\centering
\small
\begin{tabular}{@{}lrrr@{}}
\toprule
Config & Masked $\Delta_{\text{neg}}$ & Unmasked $\Delta_{\text{neg}}$ & $\Delta$ \\
\midrule
Greedy & 0.250 & 0.250 & 0.000 \\
Temperature 0.7 & 0.185 & 0.182 & -0.003 \\
Temperature 1.0 & 0.193 & 0.193 & 0.000 \\
Temperature 1.3 & 0.133 & 0.137 & 0.003 \\
Top-k 20 & 0.205 & 0.202 & -0.003 \\
Top-k 50 & 0.163 & 0.162 & -0.002 \\
Top-k 100 & 0.202 & 0.197 & -0.005 \\
Top-p 0.8 & 0.208 & 0.200 & -0.008 \\
Top-p 0.9 & 0.152 & 0.147 & -0.005 \\
Top-p 0.95 & 0.185 & 0.180 & -0.005 \\
\bottomrule
\end{tabular}
\caption{Masking-sensitivity key trace for \texttt{description / Black man vs White woman}. The sign stays positive in all settings.}
\end{table}

\subsection*{Anti-Repetition}
The anti-repetition ablation enables \texttt{no\_repeat\_ngram\_size = 3} across the full decoding
grid and compares the resulting outputs against the baseline sweep.

\begin{table}[t]
\centering
\small
\begin{tabular}{@{}lrrrr@{}}
\toprule
Config & Base $\Delta_{\text{neg}}$ & Anti $\Delta_{\text{neg}}$ & $\Delta$ gap & Sign flip \\
\midrule
Greedy & 0.250 & 0.250 & 0.000 & no \\
Temperature 0.7 & 0.185 & 0.188 & 0.003 & no \\
Temperature 1.0 & 0.193 & 0.212 & 0.018 & no \\
Temperature 1.3 & 0.133 & 0.133 & 0.000 & no \\
Top-k 20 & 0.205 & 0.185 & -0.020 & no \\
Top-k 50 & 0.163 & 0.212 & 0.048 & no \\
Top-k 100 & 0.202 & 0.198 & -0.003 & no \\
Top-p 0.8 & 0.208 & 0.160 & -0.048 & no \\
Top-p 0.9 & 0.152 & 0.205 & 0.053 & no \\
Top-p 0.95 & 0.185 & 0.195 & 0.010 & no \\
\bottomrule
\end{tabular}
\caption{Anti-repetition does not flip the sign of the highlighted Milestone 3 bias trace in any decoding configuration.}
\end{table}

Anti-repetition improved repeated 3-gram rate in 8 of 10 configurations and improved distinct-2 in
all 10 configurations. The largest quality gain occurred at temperature 0.7, where repeated 3-gram
rate fell by about 0.015 and distinct-2 rose by about 0.010. However, the main interpretive point
is again conservative: quality becomes better behaved, but the highlighted bias conclusion remains
stable. The final integrated summary stored in \texttt{outputs/reports/week5\_final\_summary.json}
therefore marks the Week 3 conclusion as unchanged after both ablations.

\section{Optional Week 4 Extension}
The proposal listed classifier replication as an optional stronger extension. The repository
includes that extension under \texttt{src/app/exai/} and \texttt{outputs/exai/}. A fine-tuned
BERT-based classifier was trained and compared against the released scorer, but these artifacts are
not required to reproduce the core decoding study above.

\begin{table}[t]
\centering
\small
\begin{tabular}{@{}lcc@{}}
\toprule
Evaluation target & Accuracy / agreement & Macro F1 \\
\midrule
Held-out test split & 0.639 & 0.483 \\
Explanation benchmark & 0.250 & 0.125 \\
Released scorer agreement on test split & 0.694 & 0.494 \\
\bottomrule
\end{tabular}
\caption{Optional Week 4 classifier-replication summary from saved \texttt{outputs/exai/eval/} artifacts.}
\end{table}

This extension is useful as a robustness and ownership exercise, but it does not change the main
claim of the project. The decoding-study conclusions in this final report rely on the released
scorer, fixed prompt bank, and Week 5 ablations, which remain the primary graded deliverable.

\section{Discussion}
The final results support a narrow but defensible claim. Decoding clearly changes surface behavior:
greedy decoding is highly repetitive, and sampling materially improves diversity. Decoding also
changes measured bias, but those shifts are smaller and less transformative than the quality
changes. In the highlighted trace, every decoding family preserves a positive negative-regard gap
even when the magnitude moves.

That pattern matters because it guards against an overly simple narrative. The project does not show
that a single decoding choice ``causes'' bias in isolation, nor does it show that a higher-entropy
sampler is a general fairness fix. Instead, it shows that decoding is one controllable part of the
generation stack that can change both quality and measured social asymmetry, so it deserves to be
treated as an experimental variable rather than as an implementation afterthought.

The ablations reinforce that interpretation. Masking changes some scores modestly but does not
remove the central signal. Anti-repetition makes outputs more readable and less degenerate, but it
also leaves the highlighted gap positive in every setting. Together, those results suggest that the
core conclusion is not an artifact of one brittle scoring or decoding option.

\section{Limitations, Ethics, and Final Deliverables}
The study remains intentionally narrow. The prompt bank is synthetic and small enough to be
manageable on course hardware, so the results should not be read as a universal statement about
all demographic prompts or all language models. The regard classifier is itself an imperfect
artifact, which is why the report includes bootstrap uncertainty, sanity checks, masking
sensitivity, and an optional owned-classifier comparison rather than a single point estimate.

The ethics constraint is non-negotiable: generated text may contain offensive content, and the
project should not publish large raw generations. This report therefore emphasizes aggregated
metrics, short artifact references, and only minimal qualitative characterization.

The final submission package now includes:
\begin{itemize}
\item this final PDF report,
\item complete references below,
\item repository code with reproducible CLI entrypoints,
\item an updated \texttt{README.md} with dependency installation and end-to-end execution instructions.
\end{itemize}

\section*{References}
\small
\begin{enumerate}[label={[}\arabic*{]}]
\item Emily Sheng, Kai-Wei Chang, Prem Natarajan, and Nanyun Peng. \textit{The Woman Worked as a Babysitter: On Biases in Language Generation}. EMNLP-IJCNLP, 2019.
\item Ari Holtzman, Jan Buys, Li Du, Maxwell Forbes, and Yejin Choi. \textit{The Curious Case of Neural Text Degeneration}. ICLR, 2020.
\item Alec Radford, Jeffrey Wu, Rewon Child, David Luan, Dario Amodei, and Ilya Sutskever. \textit{Language Models are Unsupervised Multitask Learners}. OpenAI technical report, 2019.
\item Jacob Devlin, Ming-Wei Chang, Kenton Lee, and Kristina Toutanova. \textit{BERT: Pre-training of Deep Bidirectional Transformers for Language Understanding}. NAACL-HLT, 2019.
\item \texttt{sasha/regardv3}. Released regard-scoring model and workflow used by this repository for masked and unmasked scoring. Accessed during project implementation in March--April 2026.
\item Hugging Face \texttt{transformers} library documentation and released model interfaces for \texttt{gpt2} and BERT-family checkpoints, used as the execution framework for generation and classification.
\end{enumerate}

\end{document}
